\section{Detailed Wallet-Heterogeneity Methods}
\label{sec:wallet-methods-appendix}

The secondary high-activity analysis begins on 15 June and recalculates status
for every UTC day using only earlier observations. The primary definition ranks
eligible wallets by gross observed USDC volume over the preceding seven complete
UTC days; eligibility requires at least five actions or activity on at least
three days. The highest one per cent are classified as high activity for the
following day. Expanding history uses all observations from 1 June through the
preceding day with the same eligibility rule. Full-window status is descriptive
only. Fixed 1--7 June and 1--14 June classifications remain audit artifacts and
do not determine final HA decisions.

Crossing daily wallet status with the within-market P99 indicator creates four
groups: large trades by past high-activity wallets, large trades by other
wallets, ordinary trades by past high-activity wallets, and ordinary trades by
other wallets. Price models retain the main controls and fixed effects. Wallet
sequences are chronologically ordered within wallet and market. Outcomes measure
whether another action is observed, whether its YES-equivalent direction is the
same, whether an opposite action occurs within 60 minutes, and whether a later
opposite execution price is favourable in YES-equivalent terms. Event-level
paired contrasts use paired $t$-tests and Wilcoxon signed-rank tests with Holm
adjustment within each family.

The secondary hypotheses are:
\begin{description}
  \item[HA1.] Past high-activity wallets exhibit greater same-direction continuation.
  \item[HA2.] Past high-activity status has a price-update or directional-price association after separating large and ordinary trades.
  \item[HA3.] Large size and past high activity have an additional interaction beyond their separate associations.
  \item[HA4.] Past high-activity wallets exhibit more favourable opposite actions or large-trade-associated reversal.
\end{description}

The whale analysis captures prior cumulative gross trading volume rather than trading
frequency. For each evaluation day from 15 June through 22 July, wallets with at
least one earlier trade are ranked by cumulative gross USDC across all 360
markets. The highest one per cent are the primary past-only whale group. No
minimum number of actions, active days, or markets is imposed. A trailing
seven-day definition, expanding top 0.5 and top two per cent, and continuous
within-day standardised $\log(1+x)$ cumulative volume are robustness checks.
The full-window classification is descriptive only.

The same size--wallet cells and price, sequence, favourable-opposite, and
reversal outcomes are re-estimated. WH1 tests a separate price association; WH2
tests subsequent activity and continuation; WH3 tests a large-by-wallet-status
interaction; and WH4 tests favourable opposite actions and short-horizon full
reversal. A favourable opposite action is conditional on observing an opposite
action and is not realised profit. Neither sequence continuation nor favourable
execution establishes a trading strategy or manipulation.
